\documentclass[10pt,a4paper]{article}

\usepackage[
    left=1.35cm,
    right=1.35cm,
    top=1.15cm,
    bottom=1.15cm
]{geometry}

\usepackage{enumitem}
\usepackage{hyperref}
\usepackage{xcolor}
\usepackage{titlesec}
\usepackage{fontspec}

\definecolor{darkblue}{HTML}{0F172A}
\definecolor{accentblue}{HTML}{2563EB}

\hypersetup{
    colorlinks=true,
    urlcolor=accentblue,
    linkcolor=darkblue
}

% =========================================================
% GENERAL
% =========================================================

\setlength{\parindent}{0pt}
\setlength{\parskip}{0pt}
\renewcommand{\familydefault}{\sfdefault}
\linespread{1.02}

\raggedbottom

% =========================================================
% SECTION STYLE
% =========================================================

\titleformat{\section}
    {\normalsize\bfseries\color{darkblue}}
    {}
    {0pt}
    {}
    [\titlerule]

\titlespacing*{\section}
    {0pt}
    {0pt}
    {0pt}

% =========================================================
% LISTS
% =========================================================

\setlist[itemize]{
    leftmargin=1.15em,
    itemsep=0pt,
    topsep=0pt,
    parsep=0pt,
    partopsep=0pt
}

\setlist[description]{
    style=multiline,
    leftmargin=3.35cm,
    labelwidth=3.35cm,
    labelsep=0.12cm,
    itemsep=0pt,
    topsep=0pt,
    parsep=0pt,
    partopsep=0pt
}

\begin{document}

% =========================================================
% HEADER
% =========================================================

\begin{center}

    {\LARGE \textbf{ANDRÉ NICOLAS SILVA}} \\[-2pt]

    {\normalsize \textbf{Analista de Sistemas | Requisitos, Integrações \& Produto}} \\[1pt]

    \footnotesize
    Brasil \quad • \quad
    \href{mailto:devandrenicolas@gmail.com}{devandrenicolas@gmail.com} \\[-1pt]

    \textbf{\href{https://portfolio-andrenicolas.vercel.app}{Portfólio}}
    \quad • \quad
    \href{https://github.com/devAndreNicolas}{GitHub}
    \quad • \quad
    \href{https://www.linkedin.com/in/devandrenicolas}{LinkedIn}

\end{center}

\vspace{-3pt}

% =========================================================
% SUMMARY
% =========================================================

\section*{RESUMO PROFISSIONAL}

\footnotesize

Profissional de tecnologia com formação em Sistemas para Internet e experiência na análise de necessidades de produto e regras de negócio em uma plataforma SaaS de privacidade, compliance e auditoria digital usada por \textbf{10+ empresas}. Participa da tradução de necessidades de usuários e áreas de produto em escopo técnico, fluxos e soluções de software, acompanhando implementação, testes, depuração e evolução em produção.

Atua na investigação de problemas entre interface, APIs, serviços e dados, com experiência em \textbf{SQL/PostgreSQL, APIs REST e internas, integrações, Angular/TypeScript, Go e sistemas orientados a eventos}. Também participou de entrevistas e validações de UX, instrumentação de eventos e funis no PostHog e análise de feedback para apoiar decisões sobre os fluxos do Compass.

Combina análise funcional e contexto de produto com conhecimento técnico para conversar com usuários e engenharia, esclarecer regras e acompanhar a solução até a entrega. O cargo exercido na MSPA é \textbf{Frontend Engineer}; esta variante apresenta competências transferíveis para posições de análise de sistemas, sem alterar o histórico profissional.
% =========================================================
% SKILLS
% =========================================================

\section*{COMPETÊNCIAS PRINCIPAIS}

\footnotesize

\begin{description}
    \item[\textbf{Análise de Sistemas \& Requisitos}]
    Análise de necessidades de produto, levantamento e refinamento de requisitos, regras de negócio, definição de escopo, fluxos de usuário, alinhamento com produto e engenharia

    \item[\textbf{Integração de Sistemas \& APIs}]
    APIs REST e internas, integrações entre serviços, eventos, Server-Sent Events (SSE), automações e depuração de fluxos entre sistemas

    \item[\textbf{Dados \& Persistência}]
    SQL, PostgreSQL, queries, persistência, migrações de banco de dados e investigação de comportamento relacionado a dados

    \item[\textbf{Qualidade \& Sustentação}]
    Debugging em produção, análise de problemas, testes automatizados, code review, CI/CD, suporte à produção e documentação técnica

    \item[\textbf{Produto, Usuários \& UX}]
    Entrevistas e validação de fluxos, personas, feedback de usuários, eventos e funis de produto no PostHog, acessibilidade e clareza de interface

    \item[\textbf{Tecnologias de Sistemas}]
    TypeScript, JavaScript, Angular, React, Go, Node.js, Java e Spring/Spring Boot em contribuições na MSPA; Cloudflare Workers e serviços de borda

    \item[\textbf{Arquitetura \& Segurança}]
    Arquitetura baseada em componentes, Hexagonal Architecture, separação de responsabilidades, autorização, validação de entrada e eventos de auditoria

    \item[\textbf{Desenvolvimento com IA}]
    IA generativa, Genkit, contexto e respostas estruturadas; SDD e harnesses com checks determinísticos e guardrails de qualidade e custo
\end{description}

% =========================================================
% EXPERIENCE
% =========================================================

\section*{EXPERIÊNCIA PROFISSIONAL}

\noindent
\textbf{MSPA — Frontend Engineer (contribuições full stack)}
\hfill
\textit{Abr 2025 -- Presente}

\textit{Plataforma SaaS para Privacidade, Compliance (LGPD) e Auditoria Digital}

\begin{itemize}

    \item
    \textbf{Análise de Sistemas \& Requisitos:}
    Participou da evolução do Compass e do Trust Center, produtos
    SaaS usados por \textbf{10+ empresas}, conectando necessidades
    de produto e usuários a regras de negócio, escopo técnico,
    fluxos, implementação, validação e sustentação em produção.

    \item
    \textbf{Produto \& Tomada de Decisões Técnicas:}
    Traduziu requisitos de negócio e necessidades de produto em
    soluções técnicas, contribuindo para escopo de features,
    regras de negócio, UX, arquitetura, estratégia de
    implementação e trade-offs técnicos.

    \item
    \textbf{Engenharia AI-Native:}
    Aplicou um workflow de desenvolvimento orientado por IA
    utilizando \textbf{Codex, Claude, Cursor, CodeRabbit, Skills,
    MCPs, Agents e SDD} para acelerar implementação, explorar
    soluções, compreender codebases, investigar problemas,
    revisar mudanças e iterar sobre funcionalidades de produto,
    mantendo consistência arquitetural e qualidade de código.

    \item
    \textbf{Engenharia de Produto com IA Generativa:}
    Integração de \textbf{Generative AI em produtos SaaS em
    produção}, desenvolvendo funcionalidades voltadas ao usuário
    com contextual prompt engineering, input sanitization e
    respostas estruturadas de IA utilizando \textbf{Genkit},
    com foco em funcionalidades de produto confiáveis e
    consistentes.

    \item
    \textbf{Engenharia Full Stack:}
    Implementou funcionalidades em produção atravessando
    frontend, backend, APIs, dados e infraestrutura, atuando
    entre diferentes camadas da aplicação de acordo com os
    requisitos de produto e problemas técnicos.

    \item
    \textbf{Stack Ampliada de Plataforma:}
    Contribui na MSPA com \textbf{Node.js, Java, Spring/Spring Boot,
    AWS e Azure/serviços cloud, containers Podman/Docker-compatíveis,
    Bash, Testcontainers, imagens Docker de produção, testes, acessibilidade,
    documentação técnica, entrega de UI e engenharia orientada à segurança},
    selecionando a camada adequada conforme os requisitos do produto.

    \item
    \textbf{Engenharia Backend \& Sistemas:}
    Contribuiu para \textbf{serviços backend e APIs em Golang},
    investigando e implementando soluções envolvendo
    concorrência, processamento assíncrono, eventos, integrações
    entre serviços, automações e problemas de produção em
    diferentes camadas da aplicação.

    \item
    \textbf{Dados \& Regras de Negócio:}
    Trabalhou com \textbf{SQL e PostgreSQL}, incluindo queries,
    persistência de dados, database migrations, regras de negócio
    e investigação de problemas relacionados a dados em
    diferentes workflows da aplicação.

    \item
    \textbf{Arquitetura de Software:}
    Participou de decisões arquiteturais e implementou soluções
    seguindo \textbf{Hexagonal Architecture}, arquitetura baseada
    em componentes, separação de responsabilidades, sistemas
    reutilizáveis e limites de sistema manuteníveis.

    \item
    \textbf{Web \& Edge Systems:}
    Desenvolveu um consent engine agnóstico utilizando
    \textbf{Web Components, Lit e Shadow DOM}, integrado com
    \textbf{Cloudflare Workers, KV, D1 e Durable Objects}
    para execução na edge e funcionalidades distribuídas.

    \item
    \textbf{Performance \& Confiabilidade:}
    Implementou melhorias de performance e SEO, investigando
    gargalos e otimizando comportamento de carregamento e
    performance da aplicação, mantendo usabilidade e
    confiabilidade em produção.

    \item
    \textbf{Entrega de Engenharia:}
    Atuou ao longo do ciclo de entrega de software, incluindo
    \textbf{CI/CD, deployment, debugging, investigação em
    produção e integração de mudanças entre diferentes camadas
    da aplicação}.

    \item
    \textbf{Coordenação de Engenharia:}
    Reorganizou o workflow de desenvolvimento durante uma fase
    crítica do projeto, melhorando clareza de escopo, organização
    de sprints e produtividade de uma \textbf{equipe de 6 pessoas},
    com coordenação direta de 4 integrantes.

\end{itemize}

% =========================================================
% OPEN SOURCE
% =========================================================

\section*{CONTRIBUIÇÕES OPEN SOURCE}

\noindent
\textbf{Stoat / Revolt Platforms LTD — Engenheiro Frontend Voluntário}
\hfill
\textit{2026 -- Presente}

\footnotesize

Contribui para o projeto open source com foco em
\textbf{UX engineering e desenvolvimento frontend}, identificando
problemas voltados aos usuários e transformando-os em melhorias
técnicas que tornam o produto mais claro e intuitivo para os
usuários.
\textbf{Pull request aprovado pelos maintainers do projeto.}

% =========================================================
% PROJECTS
% =========================================================

\section*{PROJETOS}

\noindent
\textbf{Marca Racha — Regras de acesso e fluxos de partidas}
\hfill
\textit{Next.js, TypeScript, Supabase, PostgreSQL, Prisma, Vitest}

\footnotesize

Produto local para organizar partidas amadoras, com partidas públicas e privadas, descoberta por região e pedidos de entrada aprovados pelo organizador. Testes comportamentais cobrem ownership, permissões por papel, membership, visibilidade e transições de estado.

\vspace{0.5pt}

\noindent
\textbf{Compass e Trust Center — SaaS de privacidade e compliance}
\hfill
\textit{Angular, TypeScript, Lit, Web Components, Cloudflare}

\footnotesize

Participação na análise de fluxos de mapeamento de dados e consentimento, entrevistas e validações de UX, eventos e funis no PostHog, além de implementação de componente de consentimento e integração com serviços de borda.

\vspace{0.5pt}

\noindent
\textbf{Terto Beats — Fluxos de pedido e pagamento}
\hfill
\textit{Next.js, TypeScript, Supabase, Stripe, Mercado Pago}

\footnotesize

Implementou criação de pedidos pendentes, validação de provedor, moeda e valor, início de checkout e tratamento de webhooks de Stripe e Mercado Pago para concluir o fluxo do pedido.

\vspace{0.5pt}

\noindent
\textbf{RendeCerto — Simulador de cenários financeiros}
\hfill
\textit{Next.js, TypeScript, Vitest}

\footnotesize

Implementou regras determinísticas de cálculo e comparação de cenários, com testes automatizados para validar resultados em diferentes entradas.

\vspace{0.5pt}

\noindent
\textbf{Saúde em Campo — Arquitetura mobile offline-first}
\hfill
\textit{React Native, Expo, SQLite, Drizzle, TanStack Query}

\footnotesize

Desenhou uma aplicação para equipes que registram visitas e informações de saúde em campo mesmo sem conectividade. A proposta combina isolamento local por usuário, fila de sincronização, idempotência, retries/backoff, RBAC e priorização de risco; é evidência de arquitetura, não de implantação clínica.

\vspace{0.5pt}

\noindent
\textbf{Missões do Dia — Aplicativo Android}
\hfill
\textit{Java, SQLite, SharedPreferences, Services, BroadcastReceiver, AsyncTask, XML}

\footnotesize

Desenvolveu um aplicativo acadêmico Android de CRUD de missões com persistência local, componentes nativos e layouts XML.

% EDUCATION
% =========================================================

\section*{FORMAÇÃO}

\noindent
\textbf{Tecnólogo em Sistemas para Internet} — UNCISAL
\hfill
\textit{Concluído em 2026}

% =========================================================
% LANGUAGES
% =========================================================

\vspace{-2pt}
\section*{IDIOMAS}

\noindent
\textbf{Português:} Nativo
\quad • \quad
\textbf{Inglês:} Intermediário/Avançado (em aprimoramento contínuo)

\end{document}
